\documentclass[../RFC-HDFG-2026-003.tex]{subfiles}

\begin{document}

\section{Test Plan}
\label{sec:test-plan}

\subsection{Blob Storage Tests}
\label{sec:test-blob}

These tests exercise \texttt{H5Pappend\_filter\_blob} and the blob
on-disk format (\SectionRef{sec:blob-format}).

\begin{itemize}
\item \texttt{blob-01} (\textbf{create/open round-trip}): Create a
  chunked dataset with a filter configured via
  \texttt{H5Pappend\_filter\_blob} using the library's default
  \texttt{H5HG}-based storage (\texttt{write\_blob}/\texttt{read\_blob}
  both \texttt{NULL}). Close and reopen the file; verify the filter's
  \texttt{read\_blob} path (default reader) returns bytes identical to
  those originally passed to \texttt{H5Pappend\_filter\_blob}.
\item \texttt{blob-02} (\textbf{custom callbacks}): Repeat blob-01 with
  a test filter that implements non-NULL
  \texttt{write\_blob}/\texttt{read\_blob}/\texttt{close\_blob}; verify
  all three are invoked at the expected lifecycle points
  (\SectionRef{sec:blob-lifecycle}) and that \texttt{close\_blob} runs
  exactly once per \texttt{H5Dclose}.
\item \texttt{blob-03} (\textbf{delete reclaims heap space}): Create a
  blob-configured dataset, note the file size, delete the dataset, and
  verify (e.g., via \texttt{h5stat} or a subsequent
  \texttt{H5Fget\_freespace} check) that the heap space is reclaimed
  rather than leaked.
\item \texttt{blob-04} (\textbf{\texttt{H5Pencode}/\texttt{H5Pdecode}
  round-trip}): Build a DCPL via \texttt{H5Pappend\_filter\_blob},
  \texttt{H5Pencode} it, and \texttt{H5Pdecode} the result into a new
  DCPL in a separate process with no access to the original file. Verify
  the decoded DCPL's blob bytes are byte-identical to the original,
  confirming \SectionRef{sec:blob-encode-decode}'s inline-serialization
  decision.
\item \texttt{blob-05} (\textbf{\texttt{h5repack} cross-file copy}): Run
  \texttt{h5repack} on a file containing a blob-configured dataset;
  verify the output file's dataset reads back correctly and that its
  blob is stored at a new, independent on-disk locator (not the source
  file's locator), per \SectionRef{sec:blob-repack}.
\item \texttt{blob-06} (\textbf{oversized default-storage blob}):
  Configure a filter via \texttt{H5Pappend\_filter\_blob} with a
  multi-megabyte buffer using the default \texttt{H5HG}-based storage;
  verify create/open succeeds and the runtime chunk I/O path is
  unaffected (compression throughput on the same file is unchanged
  relative to an equivalent dataset without a blob).
\item \texttt{blob-E01} (\textbf{malformed call}):
  \texttt{H5Pappend\_filter\_blob} with \texttt{buf=NULL} and
  \texttt{size>0}, or \texttt{buf!=NULL} and \texttt{size=0}, returns
  \texttt{H5E\_BADVALUE} and does not modify the DCPL.
\item \texttt{blob-E02} (\textbf{unknown filter}):
  \texttt{H5Pappend\_filter\_blob} with an unregistered filter
  \texttt{id} returns \texttt{H5E\_NOFILTER}.
\end{itemize}

\subsection{Parallel HDF5 Tests}
\label{sec:test-parallel}

This test requires a Parallel HDF5 build and at least two MPI ranks.

\begin{itemize}
\item \texttt{par-blob} (\textbf{rank-0 write and broadcast}): All ranks
  call \texttt{H5Pappend\_filter\_blob} with identical arguments,
  followed by a collective \texttt{H5Dcreate}. Verify that exactly one
  heap object is created (not one per rank) and that every rank's copy
  of the pipeline message carries the same locator, confirming the
  rank-0-write-and-broadcast protocol of \SectionRef{sec:blob-parallel}.
\end{itemize}

\end{document}
